\chapter{Combinatorial contracts and modulus arithmetic}

\section{Discrete renormalization data}

\begin{definition}[Renormalization profile]
\label{def:renormalization-profile}
A discrete renormalization profile consists of:
\[
 (\tau,p_n,\rho_n)_{n\in\N},
\]
where $\tau$ is one of the primitive, satellite, or neutral types,
$p_n\in\N$ is the relative period at level $n$, and
$\rho_n\in\N$ records the combinatorial rotation data.
\end{definition}

\begin{definition}[Bounded type]
\label{def:bounded-type}
The rotation profile $(\rho_n)$ is of bounded type if
\[
 \exists B\in\N\ \forall n,\qquad \rho_n\le B.
\]
\end{definition}

The Lean definitions make bounded type an explicit predicate:
\begin{itemize}
\item \lean{RenormalizationCombinatorics};
\item \lean{BoundedType};
\item \lean{IsBoundedTypeSatelliteProfile}.
\end{itemize}
They do not identify bounded type with a numerical picture or with bounded
periods: the rotation sequence is a separate field.

\begin{definition}[Cumulative depth]
\label{def:cumulative-depth}
For a profile with periods $(p_n)$ define
\[
 d_0:=0,\qquad d_{n+1}:=d_n+p_n.
\]
\end{definition}

\begin{theorem}[Depth monotonicity]
\label{thm:depth-monotonicity}
If $p_n>0$ for every $n$, then $d_n$ is monotone.  For the period-doubling
profile $p_n=2$, one has $d_n=2n$.
\end{theorem}

The proofs are elementary arithmetic and are available as
\lean{cumulativeDepth_monotone} and
\lean{cumulativeDepth_periodDoubling}.  This separates the discrete
renormalization bookkeeping from the analytic construction of the tower.

\section{Modulus divergence}

\begin{lemma}[Uniform lower bound]
\label{lem:combinatorial-modulus-lower}
Let $(a_n)$ be a sequence of nonnegative real numbers.  If
\[
 \exists\mu>0\ \forall n,\qquad \mu\le a_n,
\]
then $\sum_n a_n$ is divergent.
\end{lemma}

\begin{proof}
The $N$-th partial sum is at least $N\mu$, so the partial sums are
unbounded.
\end{proof}

The exact Lean theorem is
\lean{not_summable_of_uniform_lower_bound}; the wrapper
\lean{not_summable_of_aPrioriBounds} applies it to
\lean{APrioriBounds}.  This is the arithmetic core shared by the principal
nest arguments in the Lyubich, Dudko--Lyubich, and Kahn--Kapiamba--Lyubich
routes.

\begin{theorem}[Tower-to-puzzle arithmetic contract]
\label{thm:tower-to-puzzle-arithmetic}
Suppose a renormalization tower has cumulative depths $d_n$, and suppose its
principal-nest annulus at level $n$ is the disjoint union of ordinary puzzle
annuli with indices in $[d_n,d_{n+1})$.  If the principal-nest moduli have
the uniform lower bound $\mu>0$, then the ordinary puzzle-annulus modulus
series is not summable.
\end{theorem}

\begin{proof}
The principal-nest moduli are finite sums of disjoint ordinary annulus
moduli.  Summability of the ordinary series would imply summability of the
principal-nest series, contradicting Lemma~\ref{lem:combinatorial-modulus-lower}.
\end{proof}

\section{Source-specific combinatorial predicates}

\begin{center}\small
\begin{tabular}{@{}P{0.29\textwidth}@{\hspace{4pt}}P{0.49\textwidth}@{\hspace{4pt}}P{0.14\textwidth}@{}}
\textbf{Class} & \textbf{Discrete contract} & \textbf{Use}\\
\hline
Bounded-type satellite &
\lean{IsBoundedTypeSatelliteProfile}: satellite type plus bounded rotation
sequence. & Dudko--Lyubich\\
Period doubling &
\lean{IsPeriodDoublingProfile}: primitive type and $p_n=2$. &
Feigenbaum\\
Bounded period doubling &
\lean{IsBoundedPeriodDoublingProfile}: period doubling plus bounded
rotation data. & Feigenbaum interface\\
Neutral &
\lean{IsNeutralProfile}: neutral tower label; bubble and rigidity fields are
separate. & Dudko--Lim--Lyubich\\
\end{tabular}
\end{center}

The parabolically bounded primitive class is intentionally not reduced to
one of these elementary predicates.  Its cusp control is an analytic
estimate, not merely a bound on the integer rotation sequence.  The source
contract is:
\begin{center}
\lean{ParabolicallyBoundedPrimitiveCertificate}.
\end{center}

\begin{openestimate}[Analytic realization of combinatorics]
\label{est:combinatorial-realization}
Construct the marked renormalization tower whose discrete profile satisfies
the relevant predicate, prove the principal-nest decomposition, and supply
the modulus bounds required by the corresponding certificate.
\end{openestimate}
